\section*{Bab \ref{ch:g6:water-underground} --- Air di Bawah Tanah}

\begin{solution}{exo:g6:water-underground:1}
\omterm{def:g6:water-underground:permeability}{Permeabel}: \omterm{def:g3:stones-pebbles-sand:sand}{pasir}, \omterm{def:g3:stones-pebbles-sand:pebble}{kerikil}, kapur. \omterm{def:g6:water-underground:permeability}{Kedap air}: lempung, granit yang tidak
retak, serpih.
\end{solution}

\begin{solution}{exo:g6:water-underground:2}
(a) \omterm{def:g6:water-underground:water-table}{muka air tanah}; (b) \omterm{def:g6:water-underground:aquifer}{akuifer}; (c) \omterm{def:g6:water-underground:spring}{mata air}; (d) \omterm{def:g6:water-underground:spring}{sumur}.
\end{solution}

\begin{solution}{exo:g6:water-underground:3}
(a) Air menguap kembali ke udara; air mengalir di permukaan masuk ke
\omterm{def:g1:landscapes-around-us:water}{sungai-sungai} kecil; air meresap ke dalam tanah (\omterm{def:g6:water-underground:infiltration}{infiltrasi}). (b) Di tempat
parkir: permukaannya yang keras hampir tidak membiarkan air meresap, jadi
air mengalir pergi. Di padang rumput, tanah dan rumput membiarkan banyak air
meresap.
\end{solution}

\begin{solution}{exo:g6:water-underground:4}
(a) Salah: lempung sangat \omterm{def:g6:water-underground:porosity}{berpori} tetapi \omterm{def:g6:water-underground:permeability}{kedap air}, karena \omterm{def:g6:water-underground:porosity}{porinya} sangat
kecil. (b) Benar. (c) Salah: \omterm{def:g6:water-underground:water-table}{muka air tanah} mengikuti bentuk daratan, lebih
landai, dan naik turun mengikuti musim. (d) Benar: kira-kira 30~\% dari
seluruh \omterm{def:g2:water-on-earth:salt-fresh}{air tawar} adalah \omterm{def:g6:water-underground:groundwater}{air tanah}, dibandingkan dengan kira-kira seperempat
dari satu persen di \omterm{def:g1:landscapes-around-us:water}{danau}.
\end{solution}

\begin{solution}{exo:g6:water-underground:5}
$10 \times 25 \div 100 = 2.5$: kamu dapat menuangkan 2.5~liter sebelum ember
itu meluap. $10 \times 22 \div 100 = 2.2$: 2.2~liter akan keluar melalui
lubang itu.
\end{solution}

\begin{solution}{exo:g6:water-underground:6}
(a) $2 \times 1 \times 1 = 2$~meter kubik, yaitu 2000~liter.
(b) $2000 \times 20 \div 100 = 400$~liter. (c) $2000 \times 18 \div 100 =
360$~liter keluar; $400 - 360 = 40$~liter tetap terperangkap.
\end{solution}

\begin{solution}{exo:g6:water-underground:7}
(a) \omterm{def:g6:water-underground:spring}{Mata air} itu berada pada kira-kira 203~m, \omterm{def:g1:landscapes-around-us:water}{sungainya} pada kira-kira
132~m. (b) \omterm{def:g6:water-underground:spring}{Mata air} itu kira-kira 70~m di atas \omterm{def:g1:landscapes-around-us:water}{sungai}: airnya mengalir
menuruni lereng masuk ke \omterm{def:g1:landscapes-around-us:water}{sungai}. (c) Di dalam lereng \omterm{def:g1:landscapes-around-us:landform}{bukit} kapur, \omterm{def:g6:water-underground:water-table}{muka air
tanah} berada di bawah permukaan tanah; \omterm{def:g6:water-underground:water-table}{muka air tanah} hanya bertemu dengan
permukaan di tempat dasar lapisan kapur, yang terletak di atas serpih yang
\omterm{def:g6:water-underground:permeability}{kedap air}, muncul di lereng.
\end{solution}

\begin{solution}{exo:g6:water-underground:8}
(a) Setelah musim \omterm{def:g4:clouds-rain-snow:precipitation}{hujan}: $401 - 336 = 65$~m; setelah musim kemarau:
$401 - 314 = 87$~m. (b) $336 - 314 = 22$~m.
\end{solution}

\begin{solution}{exo:g6:water-underground:9}
Lebih dalam. \omterm{def:g6:water-underground:water-table}{Muka air tanah} mengikuti bentuk daratan, tetapi lebih landai:
di puncak \omterm{def:g1:landscapes-around-us:landform}{bukit} \omterm{def:g6:water-underground:water-table}{muka air tanah} lebih tinggi daripada di bawah \omterm{def:g6:water-underground:spring}{sumur}~A,
tetapi permukaan tanahnya naik jauh lebih banyak, sehingga airnya terletak
lebih jauh di bawah permukaan (pada penampang, lebih dari 110~m di bawah).
\end{solution}

\begin{solution}{exo:g6:water-underground:10}
\omterm{def:g6:water-underground:porosity}{Pori-pori} lempung sangat kecil: air tetap menempel pada keping-kepingnya
dan tidak dapat mengalir melalui lempung menuju \omterm{def:g6:water-underground:spring}{sumur}. Lempung itu \omterm{def:g6:water-underground:porosity}{berpori}
tetapi \omterm{def:g6:water-underground:permeability}{kedap air}, jadi hanya sedikit air (20~liter dari 500) yang keluar.
\end{solution}

\begin{solution}{exo:g6:water-underground:11}
(a) $30.1 \div 0.26 \approx 116$: lebih dari seratus kali lebih banyak.
(b) $100 - 69 = 31$: kira-kira 31~\% yang tersisa, dan \omterm{def:g6:water-underground:groundwater}{air tanah} saja sudah
kira-kira 30 di antaranya. (c) \omterm{def:g6:water-underground:groundwater}{Air tanah} adalah simpanan \omterm{def:g2:water-on-earth:salt-fresh}{air tawar} cair yang
jauh paling besar; es itu membeku dan sebagian besar letaknya jauh, di
daerah kutub dan di \omterm{def:g4:world-physical-map:range}{pegunungan}. Selain itu, \omterm{def:g6:water-underground:groundwater}{air tanah} terdapat hampir di
mana-mana di bawah kaki kita.
\end{solution}

\begin{solution}{exo:g6:water-underground:12}
Kedua \omterm{def:g6:water-underground:spring}{sumur} itu mengambil air dari \omterm{def:g6:water-underground:aquifer}{akuifer} kapur yang sama. Pompa yang kuat
itu mengeluarkan air lebih cepat daripada \omterm{def:g4:clouds-rain-snow:precipitation}{hujan} mengisi kembali kapurnya,
jadi \omterm{def:g6:water-underground:water-table}{muka air tanah} turun di seluruh sekitarnya, sampai di bawah dasar \omterm{def:g6:water-underground:spring}{sumur}
kedua, yang lalu mengering. Kedua petani itu dapat bersepakat untuk memompa
lebih sedikit, terutama pada musim panas, berbagi satu \omterm{def:g6:water-underground:spring}{sumur}, atau
memperdalam \omterm{def:g6:water-underground:spring}{sumur} kedua; yang terpenting, mereka tidak boleh mengambil air
lebih banyak daripada yang dikembalikan \omterm{def:g4:clouds-rain-snow:precipitation}{hujan}.
\end{solution}

\begin{solution}{pb:g6:water-underground:1}
\textbf{1.} Pertanian itu berdiri di atas kapur; di bawahnya terdapat
serpih. Kapur itu \omterm{def:g6:water-underground:permeability}{permeabel}, serpih tidak.

\textbf{2.} \omterm{def:g4:clouds-rain-snow:precipitation}{Hujan} meresap ke dalam kapur dan berhenti di atas serpih. Di
tempat dasar lapisan kapur muncul di lereng, \omterm{def:g6:water-underground:water-table}{muka air tanah} bertemu dengan
permukaan tanah dan airnya keluar: itulah \omterm{def:g6:water-underground:spring}{mata air}.

\textbf{3.} $401 - 336 = 65$~m di bawah permukaan tanah.

\textbf{4.} Dasar \omterm{def:g6:water-underground:spring}{sumur} itu berada pada $401 - 100 = 301$~m, di bawah \omterm{def:g6:water-underground:water-table}{muka
air tanah} yang berada pada 314~m: \omterm{def:g6:water-underground:spring}{sumur} itu masih mencapai air, sedalam $314 - 301 =
13$~m.

\textbf{5.} $100 \times 50 \times 2 = 10\,000$~meter kubik.

\textbf{6.} $10\,000 \times 20 \div 100 = 2000$~meter kubik.

\textbf{7.} $10\,000 \times 18 \div 100 = 1800$~meter kubik keluar;
$2000 - 1800 = 200$~meter kubik tetap terperangkap.

\textbf{8.} $1800 \times 1000 = 1\,800\,000$~liter.

\textbf{9.} $18 \times 1000 = 18\,000$~liter setiap hari.

\textbf{10.} $18 \times 7 = 126$~meter kubik seminggu.

\textbf{11.} \omterm{def:g4:clouds-rain-snow:precipitation}{Hujan} pada musim \omterm{def:g4:clouds-rain-snow:precipitation}{hujan} meresap ke dalam kapur dan mengisinya
kembali setiap tahun (asalkan pertanian-pertanian tidak memompa lebih banyak
daripada yang dikembalikan \omterm{def:g4:clouds-rain-snow:precipitation}{hujan}).

\textbf{12.} $1800 \div 18 = $~\textbf{100 hari}: sedikit lebih dari tiga
bulan tanpa \omterm{def:g4:clouds-rain-snow:precipitation}{hujan} sama sekali.
\end{solution}
